\documentclass[../../main.tex]{subfiles}
\begin{document}

\begin{center}
  {\fontsize{48pt}{58pt}\selectfont \textbf{1965}\normalsize\textbf{年度}}
\end{center}

\begin{table}[htb]
  \centering
  \caption{出題テーマと難易度}
  \begin{tabular}{|c|p{6.2cm}|c|c|c|}
    \hline
    \textbf{問題番号} & \textbf{テーマ} & \textbf{思考量} & \textbf{計算量} & \textbf{総合} \\
    \hline
    第1問 & 円周上の点と直線の交点をめぐる対話 & 易 & 易 & 難 \\
    \hline
    第2問 & 定数根をもつ二次方程式 & 易 & 易 & 易 \\
    \hline
    第3問 & 三角関数の恒等式と角の決定 & 易 & 易 & 易 \\
    \hline
    第4問 & 角の二等分線上の点の判定 & 難 & 難 & 難 \\
    \hline
    第5問 & 微分係数の定義と計算 & 易 & 易 & 易 \\
    \hline
    第6問 & 積分方程式から多項式を決定 & 易 & 易 & 易 \\
    \hline
  \end{tabular}
\end{table}

\begin{table}[htb]
  \centering
  \caption{解答}
  \renewcommand{\arraystretch}{1.4}
  \begin{tabular}{|c|c|p{9cm}|}
    \hline
    \textbf{問題番号} & \textbf{小問} & \textbf{答え} \\
    \hline
    第1問 & -- &  \\
    \hline
    第2問 & -- & $a=2,b=8$．他方の根は$k=2$で最小値$6$ \\
    \hline
    第3問 & -- & $(\alpha,\beta)=(\pm\dfrac{\pi}{6},\mp\dfrac{\pi}{6})$（複号同順） \\
    \hline
    第4問 & -- & $A$は角$XOY$の二等分線上にある（証明問題） \\
    \hline
    第5問 & (ア) & $f\prime(1)=-3$ \\
    \cline{2-3}
    & (イ) & $f\prime(1)=\dfrac{\sqrt{3}}{2}$ \\
    \hline
    第6問 & -- & $f(t)=12t^2-4$ \\
    \hline
  \end{tabular}
\end{table}

\paragraph{試験形態}
6問．試験時間は要確認．

\end{document}
